% ======================================================================
% Section 3 — Problem Formulation and Graph Learning Framework
% Compatible with \documentclass[pdflatex,sn-mathphys-num]{sn-jnl}
% Citations: natbib \citep/\citet against MyReferences.bib
% ======================================================================
\section{Problem Formulation and Graph Learning Framework}\label{sec:method}

This section defines the forecasting task, the graph constructions, and the
mechanisms by which the forecasting backbones consume graph structure. Two
design dimensions are kept strictly separate throughout: the \emph{source} of
the graph (physical, identity, learned contemporaneous, or learned multi-lag)
and the \emph{consumption} of the graph (one static operator versus
timestep-specific operators). Throughout, $\mathbf{A}$ denotes the binary
adjacency consumed by the forecasting backbone, whereas $W$ denotes the
continuous coefficient matrix estimated by the graph-structure-learning
procedure; the two symbols are not interchangeable.

\subsection{Forecasting Task and Model Baselines}\label{sec:task}

Let $x(t)\in\mathbb{R}^{N}$ denote the traffic speeds observed at $N$ sensors
at time step $t$, where one step corresponds to the sampling interval of the
dataset (5\,min for Los-loop, 15\,min for SZ-Taxi; Sect.~\ref{sec:setup}).
Given an input window of length
\[
T=12,
\]
the forecasting task is to map
\[
\underbrace{\bigl[x(t-11),\,\dots,\,x(t)\bigr]}_{\displaystyle X\in\mathbb{R}^{T\times N}}
\qquad\longmapsto\qquad
\underbrace{\bigl[x(t+1),\,\dots,\,x(t+\mathrm{PH})\bigr]}_{\displaystyle Y\in\mathbb{R}^{\mathrm{PH}\times N}},
\]
where the prediction horizon $\mathrm{PH}\in\{1,2,3,4\}$ is expressed in steps
of the sampling interval. The $\mathrm{PH}$-dimensional target at every sensor
is produced jointly in a single forward pass by a shared linear output layer;
prediction is multi-step but \emph{not} autoregressive, and no predicted value
is fed back into the model.

The backbone adjacency is a binary matrix
\[
\mathbf{A}\in\{0,1\}^{N\times N}.
\]
It is produced from a continuous coefficient matrix
\[
W\in\mathbb{R}^{d\times d},
\qquad
d=\begin{cases}
N & \text{(contemporaneous construction, Sect.~\ref{sec:gsl_contemp})},\\
(L+1)N & \text{(multi-lag construction, Sect.~\ref{sec:gsl_multilag})}.
\end{cases}
\]
$\mathbf{A}$ denotes the backbone adjacency, whereas $W$ denotes the continuous
DAGMA coefficient matrix.

Two reference configurations fix the ends of the graph-source spectrum:

\begin{itemize}
\item \textbf{Graph-free control.} $\mathbf{A}=\mathbf{I}_N$, the identity
adjacency, named \emph{T-GCN-NoSpatial} and \emph{GCN-NoSpatial} in the two
backbone families. The identity adjacency is passed through exactly the same
graph-normalization pipeline as every other graph configuration
(see ``Graph normalization'' below), so this control isolates the contribution of
spatial structure rather than of the graph operator itself.
\item \textbf{Physical baseline.} The dataset-provided road-network adjacency,
named \emph{T-GCN} in the T-GCN family and \emph{GCN} in the GCN family.
\end{itemize}

The two backbones are deliberately lightweight. GCN applies graph convolution
to the flattened input representation and produces the hidden state in a single
step; T-GCN incorporates the same graph convolution into a GRU-based recurrent
architecture that updates a hidden state at every input step. Generic
derivations of both architectures are given in
Sect.~2; here we only note the properties that matter for
the experimental design: all graph configurations are consumed by
architecturally identical backbones, differing only in the adjacency supplied.

\paragraph{Graph normalization.}
Every adjacency --- physical, identity, and all learned graphs --- is
transformed by the same operator, following the GCN convention of
\citet{kipf2017semi}. With $\tilde{\mathbf{A}}=\mathbf{A}+\mathbf{I}$ and
$\tilde{\mathbf{D}}_{ii}=\sum_j \tilde{\mathbf{A}}_{ij}$, the operator is
\begin{equation}\label{eq:norm}
\hat{\mathbf{A}}
=\tilde{\mathbf{D}}^{-1/2}\,\tilde{\mathbf{A}}^{\top}\,\tilde{\mathbf{D}}^{-1/2}.
\end{equation}
Self-loops are therefore added before normalization, the normalization is
symmetric, and the transposed form in Eq.~\eqref{eq:norm} is the one actually
implemented. Because $\tilde{\mathbf{A}}$ is symmetric whenever $\mathbf{A}$ is
(as holds for the physical adjacency after symmetrization-free loading of its
symmetric datasets, the identity, and the cGSL graphs), the operator coincides
with the standard symmetrically-normalized adjacency in those cases; for the
directed learned graphs it realizes the same row/column scaling on the
transposed matrix. For the identity adjacency the operator reduces exactly to
$\hat{\mathbf{A}}=\mathbf{I}$, confirming that the graph-free control differs
from the other configurations only in the absence of off-diagonal structure.

\subsection{Contemporaneous Graph Learning (GSL and cGSL)}\label{sec:gsl_contemp}

The first learned construction estimates a single dependency graph among the
$N$ sensors from \emph{simultaneous} observations. We use DAGMA
\citep{Bello2024DAGMA}, which estimates a weighted adjacency $W\in\mathbb{R}^{N\times N}$
by minimizing a linear structural-equation-model score --- an $\ell_2$ loss on
the reconstructed observations --- augmented with an $\ell_1$ sparsity penalty
$\lambda_1\lVert W\rVert_1$, subject to the DAGMA acyclicity characterization
\[
h(W) = -\log\det\!\bigl(d\mathbf{I} - W\circ W\bigr) + d \;\geq\; 0,
\]
which holds if and only if $W\circ W$ is an M-matrix, i.e., $W$ is acyclic;
$h$ is enforced numerically during the constrained optimization.

\paragraph{Exact construction.}
The DAGMA input is the subsampled matrix of training-split snapshots
\[
\mathbf{X} = \mathrm{train\_norm}[\,0::\mathrm{PH}\,],
\]
that is, every $\mathrm{PH}$-th training row, so each row is one simultaneous
observation of all $N$ sensors and columns correspond to sensors. There are
$N$ variables and \emph{no} lag blocks. For Los-loop the inputs contain
$1612$, $806$, $538$, and $403$ rows for $\mathrm{PH}=1,\dots,4$; for SZ-Taxi
the corresponding counts are $2380$, $1190$, $794$, and $595$. The
construction is therefore PH-dependent through its input subsampling, in
contrast to the multi-lag construction of Sect.~\ref{sec:gsl_multilag}.

The sparsity penalty is $\lambda_1=0.02$ for Los-loop and $\lambda_1=0.01$ for
SZ-Taxi, and DAGMA's internal magnitude threshold is set to $0.3$: after the
fit, coefficients with $|W_{ij}|<0.3$ are zeroed inside the estimation
procedure. The binary support is then extracted with the absolute-magnitude
rule
\[
\mathbf{A}_{\mathrm{GSL}} = \mathbf{1}\!\left(|W|>0\right),
\]
with the diagonal removed. Under this construction the support rule is
equivalent to retaining exactly those coefficients whose magnitude reaches the
internal threshold of $0.3$. This $0.3$-based regime is specific to the
contemporaneous construction and must not be conflated with the $0.1$
consumer threshold used by the multi-lag construction
(Sect.~\ref{sec:gsl_multilag}). The graph is learned using training data only.

The resulting directed off-diagonal edge counts are $28$ for Los-loop and $8$
for SZ-Taxi, identical across $\mathrm{PH}=1,\dots,4$.

\paragraph{cGSL.}
The second learned configuration is not a separately learned graph: it is the
symmetrized version of the very same GSL artifact,
\[
\mathbf{A}_{\mathrm{cGSL}}
=\mathbf{1}\!\left(\mathbf{A}_{\mathrm{GSL}}+\mathbf{A}_{\mathrm{GSL}}^{\top}>0\right),
\]
with the diagonal removed. No second DAGMA fit and no second threshold are
involved. The symmetrized graphs contain $56$ (Los-loop) and $16$ (SZ-Taxi)
directed entries. The same graph artifacts are shared by the GCN and T-GCN
counterparts, so both backbone families see identical adjacencies.

\paragraph{Interpretation restriction.}
The contemporaneous DAGMA graph is a statistical dependency structure
estimated from simultaneous observations. It is \emph{not} a temporal graph:
an entry $\mathbf{A}_{ij}=1$ does not mean that sensor $j$ at time $t$ predicts
sensor $i$ at time $t+1$. Acyclicity is an optimization and fitting constraint
of the estimator and is not treated as evidence of a causal structure.

\subsection{Multi-Lag Graph Learning}\label{sec:gsl_multilag}

The second learned construction estimates dependency structure between lagged
and current observations. With lag depth
\[
L=3,
\]
the lag-stacked representation is
\[
Z=\bigl[\,x(t-3),\;x(t-2),\;x(t-1),\;x(t)\,\bigr],
\]
a matrix with $(L+1)N$ columns: $4\times 207=828$ variables for Los-loop and
$4\times 156=624$ for SZ-Taxi. The canonical construction contains $1609$
training windows for Los-loop and $2377$ for SZ-Taxi. DAGMA is fitted
\emph{once} on the training split, and the learned artifact is shared across
$\mathrm{PH}=1,\dots,4$:
\begin{quote}
The multi-lag graph is PH-independent: the same learned graph artifact is
consumed for all four prediction horizons. Unlike the multi-lag construction,
the contemporaneous graph is PH-dependent because its input is subsampled as
$\mathrm{train\_norm}[0::\mathrm{PH}]$.
\end{quote}

\paragraph{Multi-lag optimization.}
The $\ell_1$ coefficient is $\lambda_1=0.01$ for both datasets. In contrast to
the contemporaneous construction, DAGMA's internal threshold is $0.0$: raw
continuous weights are retained, and thresholding is deferred to the consumer.
The two threshold regimes are kept explicit throughout:
\begin{enumerate}
\item \emph{Contemporaneous GSL:} internal DAGMA threshold $0.3$; the support
is extracted after that internal threshold.
\item \emph{Multi-lag:} internal DAGMA threshold $0.0$; raw weights retained;
the consumer thresholds at $|W|>0.1$, with the diagonal removed.
\end{enumerate}

\paragraph{Block interpretation.}
In the DAGMA convention $W_{ij}$ quantifies the estimated dependency of
variable $j$ on variable $i$. The block
$W[\,lN:(l{+}1)N,\;LN:(L{+}1)N\,]$ therefore collects, for each pair of
sensors, the estimated dependency of the current snapshot on the snapshot
$L-l$ steps earlier; after the consumer threshold this yields the lag-specific
adjacency $\mathbf{A}_k$, $k\in\{1,2,3\}$, read as ``activity at $t-k$
associated with activity at $t$''. The contemporaneous block
($x(t)\rightarrow x(t)$) is present in the DAGMA fit but is \emph{not}
consumed by any of the forecasting methods in this study.

\paragraph{Edge counts.}
After the consumer threshold, the per-lag off-diagonal edge counts are:

\begin{itemize}
\item \textbf{Los-loop:} lag-1: $12$; lag-2: $3$; lag-3: $15$; total lagged
edge slots: $30$; union: $28$ distinct off-diagonal edges.
\item \textbf{SZ-Taxi:} lag-1: $0$; lag-2: $0$; lag-3: $2$; total: $2$; union: $2$.
\end{itemize}
Edge \emph{slots} count entries per lag block and may repeat an edge across
lags; the \emph{union} counts distinct sensor pairs after merging the three
blocks.

\paragraph{Interpretation restrictions.}
The learned lag graphs are descriptive statistical dependency estimates. The
lag-structured reading is consistent with their construction and with the
consumption experiments below, but the graphs are not independently validated
as causal or predictive structures, and DAGMA's acyclicity constraint is not
causal evidence. The graphs are static: they are fitted once on the training
split and never change across the input window or across evaluation samples;
only the manner in which they are \emph{consumed} can vary with the input
timestep, as described next.

\subsection{Graph Consumption Mechanisms}\label{sec:consumption}

Graph source and graph consumption are treated as separate design dimensions.
A graph with a given edge set can be consumed in different ways, and the same
consumption mechanism can be applied to different graphs. This subsection
defines the four consumption mechanisms used in the study; the corresponding
empirical comparison is presented in Sect.~5.

\paragraph{Static graph mechanisms.}
For the physical baseline, the graph-free control, and the two contemporaneous
learned graphs (GSL and cGSL), one adjacency is used for all input steps of
the window. Graph-side extra parameters: $0$.

\paragraph{Fixed cyclic lag assignment (T-GCN-MultiGSL).}
The three lag-specific graphs $\mathbf{A}_1,\mathbf{A}_2,\mathbf{A}_3$ are
consumed separately, one per input step. At input step $t$, where $t=0$ denotes
the most recent input step of the window, the graph index follows the rule
\begin{equation}\label{eq:lagidx}
\mathrm{graph\_idx} = (T-1-t) \bmod 3,
\end{equation}
so the most recent input step is convolved with the lag-1 graph, the preceding
step with the lag-2 graph, and the oldest step of the window with the lag-3
graph; the assignment cycles for the remaining steps of the 12-step window.
Extra parameters: $0$.

\paragraph{Global lag mixture (T-GCN-MultiGSL-Weighted).}
Three learned scalar weights, one per lag graph, are passed through a softmax,
and the resulting convex combination of the three normalized lag-graph
operators is used for every input step. Extra parameters: $+3$. This variant
is treated as a supplementary mechanism.

\paragraph{Per-node, per-timestep gate (T-GCN-MultiGSL-Mix).}
At each input step $t$ and each node, a small gate network takes the pair
$[x_t;\,h_{t-1}]$ --- the current input and the recurrent hidden state --- and
produces a softmax over the $K=3$ lag graphs. The gate output mixes the three
normalized lag-graph operators into one operator for that node and step,
which is then applied before the GRU-style update. Extra parameters: $+4{,}419$.
Crucially, Mix varies \emph{how} the static lag-specific graphs are consumed;
it does not relearn or modify the graph edge set at any timestep, and it is
not an adaptive-graph method: the learned structure itself never changes.

\paragraph{Static union in GCN (GCN-MultiGSL).}
The GCN backbone consumes the whole input window in a single graph
convolution, so per-timestep graph assignment is not applicable. GCN-MultiGSL
therefore merges the three lag graphs into one static adjacency for the
complete window,
\[
\mathbf{A}_{\cup} = \bigvee_{k=1}^{3} \mathbf{A}_k,
\]
the element-wise maximum for binary adjacency matrices, normalized identically
to all other graphs. The same multi-lag graph artifacts are thus consumed by
GCN-MultiGSL through a static union and by T-GCN-MultiGSL through
lag-specific, per-timestep consumption. The comparison is designed to assess
whether the utility of the same learned multi-lag structure depends on how it
is consumed; because GCN and T-GCN also differ in backbone architecture, the
comparison does not isolate consumption as the sole differing factor.

\paragraph{Parameter counts.}
All configurations within a backbone family share the same backbone capacity,
and the graph source contributes no trainable parameters. Verified from the
implementation, the trainable parameters (excluding the shared linear output
layer) are $768$ for GCN and $12{,}672$ for T-GCN in every single-graph
configuration; T-GCN-MultiGSL-Weighted adds $3$ and T-GCN-MultiGSL-Mix adds
$4{,}419$ parameters relative to T-GCN-MultiGSL. The capacity implications of
the Mix gate are examined separately in the experiments
(Sect.~\ref{sec:setup}).
